% ---
% id: protocol-compiler/def:spec
% kind: def
% title: "Protocol specification"
% summary: >-
%   A protocol specification is the compiler's input: the transmitter
%   artefact, the receiver part, a two-segment reduction chain, the public part
%   and the constraint entries. In the general format the receiver part and the
%   chain come per measured factor, with a transcript and scoped entries.
% prominence: standard
% status: implemented
% uses: [def:homog, def:declaration, transmitter/def:source-spec, transmitter/def:srcert, transmitter/def:source-segment]
% mentions: [def:weighted, def:validator, def:chaincond, def:eligible, reductions/def:chain, reductions/def:link, reductions/thm:emitted-product, req:pass-through, req:derived-objects, rem:gate-by-certificate-form, rem:chain-and-validator, rem:party-label-restriction, reductions/req:factorised-emission, transmitter/ass:ens]
% assumes: []
% refs: []
% provenance: newly-derived
% implements: certiqs_key.protocol_compiler.spec:ProtocolSpecification
% lean:
% notes: "reductions/def:chain and reductions/def:link are mentions and not uses: the record is typed in this module's words, and a uses edge would carry the producer's definitional ledger into every compiler theorem, the Band-A theorem included. The transmitter edges of (a) are uses: the source part is the transmitter's definition and carries no ledger of its own. def:chaincond is a mention: it points forward to the conditions decided on the record."
% ---

The compiler's input comes in five parts. Parts (a)--(e) describe the
specification of a declaration in the \tref{basic format}; the paragraphs on
the general format read the same five parts more widely. The descriptors a
link carries are, for a \tref{source link}, the transmitter's
\tref{source descriptor} and \tref{target descriptor}; for a receiver link
they are receiver-side. How the compiler treats the slots it reads is
Requirement~\ref{req:pass-through}. Which objects check V1 compares, and why
a supplied witness is compared at check V7 instead, is
Remark~\ref{rem:chain-and-validator}.

In the general format, how (DX1) restricts the \tref{party partition} is
Remark~\ref{rem:party-label-restriction}; that the product rule
\eqref{eq:productfamily} of (b$'$) is formed by no one is
Requirement~\ref{req:derived-objects}; and the ordered
list of the factors' artefacts in (c$'$) is the one
\xref{reductions/req:factorised-emission} emits. A certificate declared in
\tref{block form} is refused at gate (G2) of Definition~\ref{def:eligible}.
The ground is not that this module demands purity of every certificate: at several Alice sectors global purity
fails for a presentation too. What separates the two forms is stated in
Remark~\ref{rem:gate-by-certificate-form}, and the operative ground for the
refusal is the scope rule: admitting a block form would widen an
unsafe-direction gate at the degenerate values.

\begin{definition}\label{def:spec}
A \term[The compiler's input: a source part, a receiver part, a reduction chain, a public part and a constraint part.]{protocol specification} is a tuple
\begin{equation*}
  \Spec=\big(\Spec_{\mathrm{src}},\ \Spec_{\mathrm{rec}},\ \pcChain,\
  \Spec_{\mathrm{pub}},\ \Spec_{\mathrm{con}}\big)
\end{equation*}
whose parts are:

\emph{(a) Source part $\Spec_{\mathrm{src}}$.} A \tref{source specification}
in the sense of \xref{transmitter/def:source-spec}, presented to this module as
the \emph{transmitter artefact} in its declared mode: a finite non-empty
alphabet $\Xset$ with its \tref{declared order}, a finite-dimensional system
$A$ with declared computational basis, a POVM $\{M_x\}_{x\in\Xset}$ on $A$
--- itself possibly presented as a \tref{choice-weighted family}
(Definition~\ref{def:weighted}) --- and the \tref{specified signal ensemble}
$\{p_x,\rho_x\}_{x\in\Xset}$ on the \tref{transmitted system} $\Asig$
(\xref{transmitter/ass:ens}). The mode is read off the artefact's
\tref{marginal slot} and is never a second datum. In
\tref{trusted-source mode} the slot holds $\tau_A\in\Seq{A}$ --- at $d_A=1$
the slot the transmitter filled with $[1]$ included --- and at $d_A\ge2$ the
artefact additionally carries a \tref{source-replacement certificate}
$(\tau_A,\Psi,\{p_x,\rvirt_x\}_{x\in\Xset})$ on $A\Avirt$ in the sense of
\xref{transmitter/def:srcert}, written against the \tref{virtual ensemble}
the \tref{source segment} \emph{starts from}, which lives on the
\tref{virtual transmitted system} $\Avirt$, the \tref{source descriptor} of
$\Xi_1$ (\xref{transmitter/def:source-segment}\,(S1), (S3)). When
$L_{\mathrm s}=0$ that system and that ensemble are the specified ones:
$\Avirt=\Asig$ and $\rvirt_x=\rho_x$. In \tref{untrusted-source mode} the slot
holds $\nomarg$ and there is no certificate.

\emph{(b) Receiver part $\Spec_{\mathrm{rec}}$.} A finite-dimensional physical
receiver system with a POVM $\{\Fph_w\}_{w\in\Vph}$, possibly presented as a
choice-weighted family; and, optionally, a public choice whose weight
depends on the private setting of the measuring apparatus --- a
stochastic matrix $P(f\,|\,t)$ with the setting map
$t:\Vset\to\mathcal{T}$ it is read through --- to be applied by
this module \emph{after} the \tref{reduction chain}, in the
\tref{setting-dependent form} of Definition~\ref{def:weighted}.

\emph{(c) Reduction chain $\pcChain=(\Chsrc,\Chrec)$.} A reduction chain in
the sense of \xref{reductions/def:chain}, presented to this module as a
record: an ordered \emph{pair} of ordered lists of \tref[typed link]{typed links}, the
\emph{source segment} $\Chsrc=(\Xi_1,\dots,\Xi_{L_{\mathrm s}})$ with
$L_{\mathrm s}\ge0$, which is a source segment in the sense of
\xref{transmitter/def:source-segment} --- its links act on the transmitted
system $\Asig$ and replace Alice's signal ensemble --- and the \emph{\tref[receiver segment]{receiver
segment}} $\Chrec=(\Lambda_1,\dots,\Lambda_L)$ with $L\ge1$, whose links act on
the receiver system and replace Bob's measurement. What this module reads of a
link, stated in its own words, is this: each link of either segment carries a
type tag, source and target system descriptors, a status tag, its channel as
data where the link type has one, the constraint entries it transports, and a
licence slot. A receiver-stage link in the sense of
\xref{reductions/def:link} carries more --- its \tref{reachable-set statement} and
its carried operator pair --- and this module reads none of it. The last link
of $\Chrec$ has target descriptor $(\Qsys,\SecSet,(d_s)_s)$ with the
sector projectors $\Pi_s$, and \tref[target povm]{target POVM} $\{F_v\}_{v\in\Vset}$. The two
segments are separate lists because they have different types; the pair is
composed source segment first. The \tref{interface conditions on a chain} that
this module decides on the record are those of
Definition~\ref{def:chaincond}.

\emph{(d) Public part $\Spec_{\mathrm{pub}}$.} A \tref{party partition}
$\Party$ (readouts $\rd{\mathsf A},\rd{\mathsf B}$ onto $\Oset{\mathsf A}$,
$\Oset{\mathsf B}$, and the party labels), \tref[announcement map]{announcement maps}
$a_A:\Oset{\mathsf A}\to\CAset$ and $a_B:\Oset{\mathsf B}\to\CBset$, a sifting
predicate $\varsigma:\Kann\to\{\mathrm{pass},\mathrm{fail}\}$, a key rule
$r_{\hold}^{0}:\Oset{\hold}\times\Kann\to\Zval$, a holder
$\hold\in\{\mathsf A,\mathsf B\}$, a finite acceptance alphabet $\Chatset$ and
a map $\phi:\Kann\to\Chatset$. All maps are total.

\emph{(e) Constraint part $\Spec_{\mathrm{con}}$.} A finite, possibly empty,
list of entries. Each $\hat O_k$ is presented as a \tref{block list} on
$\Qsys$ together with the \tref{sector map} it is written against, which is
carried with it as for every other supplied $\Qsys$-operator family, and is
Hermitian block by block. Each entry carries a right-hand-side datum of one of
two types, and the type fixes the entry's tag: \tref{homogeneous} with a real
$c_k$, or \tref{general} with a Hermitian $\Theta_k$ on $A$, in trusted-source
mode; a real $c_k$ in untrusted-source mode (Definition~\ref{def:homog}). The
constraint part also carries the optional attack-set \tref{licence record} and an
optional feasibility witness $\omega^{\star}$, likewise a block list with its
own map. The map each $\hat O_k$ carries is compared with the emitted $\Smap$
by check V1, an entry declared here being on the same footing as one
transported by a link. The witness's map is not compared there: $\omega^{\star}$
is a block list on $A\Qsys$ and not a $\Qsys$-operator, it is no slot of
\eqref{eq:declaration}, and branch~(c) of check V7 is where it is required to
be a block list against $\Smap$ (Definition~\ref{def:validator}).

\emph{The general format.} A specification to be compiled into the general
format of Definition~\ref{def:declaration}\,(DX1)--(DX6) has the same five
parts, read as follows. At one Alice sector, one \tref{measured factor},
\tref{POVM form}, two messages neither of which reads the other and no
party-local entry these readings are (a)--(e) above, part for part on every
component a derivation reads: at those values every such object named below
is the object of the same name there and nothing else is presented. On the
party labels of the tensor factors of each $\Hs{\Qsys_s}$, which no
derivation of this module reads and which the consuming engine's laboratory
premise reads as naming the parties' laboratories, (DX1) restricts the party
partition of (d). The party set is unchanged: it is
$\{\mathsf A,\mathsf B\}$, with the two \tref[readout map]{readout maps} of $\Party$. What
is new is that several measured factors may carry a party label and that the
factor tagged $A$ carries one too (DX1).

\emph{(a$'$) Source part.} The artefact's marginal slot may hold a block list
$(\tau^{(t)}_A)_{t\in\mathbb{T}}$ along the Alice sectors the artefact
declares, and Alice's measurement arrives in one of the two forms of (DX4):
the POVM slot, or, in trusted-source mode only, the ensemble slot holding the
block lists $\{N_x\}_{x\in\Xset}$ of \eqref{eq:ensembleform}. The
source-replacement certificate may accordingly be declared in
\tref{pure form}, in \tref{block form}, or as an
\tref{ensemble presentation}, all three in the sense of
\xref{transmitter/def:srcert}. Which form a certificate is in is read off the
artefact and is never a second datum, exactly as the mode is. That a form may
be \emph{presented} is the transmitter's rule and is stated here. Which forms
this module \emph{admits}, and what gate (G2) decides of each form, is
Definition~\ref{def:eligible}'s and is stated once there.

\emph{(b$'$) Receiver part.} $\Spec_{\mathrm{rec}}$ is an ordered list, one
entry per measured factor $f\in\mathbb{F}$ and in the order of $\mathbb{F}$: a
party label; a finite-dimensional physical receiver system with a POVM
$\{F^{\mathrm{ph},f}_{w_f}\}_{w_f\in\mathcal{V}^{\mathrm{ph},f}}$, possibly
presented as a choice-weighted family; and, optionally, a
setting-dependent public choice at that factor --- a stochastic matrix
$P_f(\,\cdot\mid\cdot\,)$ with the setting map on that factor's alphabet it is
read through --- to be applied after that factor's segment, in the
setting-dependent form of Definition~\ref{def:weighted}. No entry of the list
carries an operator on another factor's system, and no product of physical
families is presented: the product is the rule \eqref{eq:productfamily},
formed by no one.

\emph{(c$'$) Reduction chain.} $\pcChain=(\Chsrc,\Chrec)$ with $\Chsrc$ as
in (c) and $\Chrec=(\pcChain^{f}_{\mathrm{rec}})_{f\in\mathbb{F}}$ one
receiver segment per measured factor, each
$\pcChain^{f}_{\mathrm{rec}}=(\Lambda^{f}_1,\dots,\Lambda^{f}_{L_f})$,
$L_f\ge1$, a list of receiver-stage links whose source descriptor is factor
$f$'s physical receiver system and whose last target descriptor is
$(\Qsys_f,\SecSet^{f},(d_{f,s_f})_{s_f})$ with target POVM
$\{F^{f}_{v_f}\}_{v_f\in\Vset^{f}}$. What this module reads of a link is what
(c) lists and nothing more, of each segment separately: the segments are
separate lists because their links act on different systems, and no link of
one is composable with a link of another. Where the factors' data are the
ordered list of the factors' artefacts that the reductions module emits,
factor $f$'s receiver segment and \tref{target family} are that factor's
own, and the product rule of \xref{reductions/thm:emitted-product} is
emitted in place of a \tref{joint family}; that list
is what the factor list of (DX1) is built from, in its order. The interface
conditions on a chain that this module decides on the record are those of
Definition~\ref{def:chaincond}, per factor for the receiver segments.

\emph{(d$'$) Public part.} In place of $a_A$ and $a_B$ the public part
declares the \tref[transcript protocol]{transcript} of (DX5): the number $L\ge1$ of messages and,
for each $m$, its sender $\pi_m$, its finite alphabet $\mathcal{C}_m$ and
its total table
$a_m:\Oset{\pi_m}\times\mathcal{C}_1\times\dots\times\mathcal{C}_{m-1}
\to\mathcal{C}_m$. The sifting predicate $\varsigma$, the key rule
$r^{0}_{\hold}$, the holder $\hold$, the acceptance alphabet $\Chatset$ and
$\phi$ are declared on the derived $\Kann=\prod_{m}\mathcal{C}_m$ of
\eqref{eq:transcript}, which is a set determined by the tables and declared by
no one. All tables are total. The party partition $\Party$ is as in (d), with
a party label per measured factor and, where there is one, for the factor
tagged $A$.

\emph{(e$'$) Constraint part.} Each entry presents, in addition to what (e)
lists, its \tref{scope}: the label of one measured factor of the chain; and
$\hat O_k$ is a block list on $\Qsys_{f_k}$ together with the sector map of
that factor, Hermitian block by block. Its right-hand-side datum is of one of
the three forms of Definition~\ref{def:homog}: a real $c_k$ tagged
homogeneous, a Hermitian $\Theta_k$ tagged general and then itself a block
list along $\mathbb{T}$, or a real $c_k$ tagged \tref{party-local}. The
arriving tag is compared with the type of that datum at check V6, which also
decides that the scope names a measured factor of the emitted map. An optional
witness $\omega^{\star}$ is a block list over the \tref[joint sector]{joint sectors} with
its own map, compared at branch~(c) of V7 and nowhere else, for the reason
(e) gives.
\end{definition}
